% ==================================================================
% id:       prob-008
% topic:    מבוא להסתברות
% subtopic: עקרון ההכלה וההפרדה
% solution: full
% ==================================================================
בקורס בהסתברות היו שני מבחנים. 80\% מהסטודנטים עברו את המבחן הראשון, 75\% עברו את המבחן השני, ו-60\% עברו את שניהם. בוחרים סטודנט באקראי.
\begin{enumerate}
\item מה ההסתברות שהוא עבר את המבחן הראשון ונכשל במבחן השני?
\hint{סמנו $A$ = „עבר את הראשון”, $B$ = „עבר את השני”. אפשר גם לסדר את הנתונים בטבלה של $2\times2$: עבר/נכשל בראשון, ועבר/נכשל בשני.}
\begin{solbox}
$A$ הוא איחוד זר של $A\cap B$ ושל $A\cap B^C$, ולכן $\PP(A\cap B^C)=\PP(A)-\PP(A\cap B)=0.8-0.6=\ans{0.2}$.
\end{solbox}
\item מה ההסתברות שהוא עבר לפחות מבחן אחד?
\begin{solbox}
לפי הכלה והפרדה: $\PP(A\cup B)=0.8+0.75-0.6=\ans{0.95}$.
\end{solbox}
\item מה ההסתברות שהוא נכשל לפחות במבחן אחד?
\begin{solbox}
„נכשל לפחות באחד” הוא המשלים של „עבר את שניהם”: $1-\PP(A\cap B)=1-0.6=\ans{0.4}$.
\end{solbox}
\end{enumerate}
